\section{串}\label{sec:Strings}\index{string|(}\thispagestyle{empty}

%<*def:Alphabet>
一个 \definend{字母表（alphabet）}\index{alphabet}
是由 \definend{符号（symbol）}\index{symbol} 或
\definend{记号（token）}\index{token} 构成的一个非空且有限的集合。
我们用打字机风格的字体书写符号，例如 \str{1} 或 \str{a}，
因为用引号将它们括起来的做法会显得笨拙。\footnote{%
  打字机字体也能将符号 ‘\protect\str{a}’
  与变量 ‘$a$’ 区分开，
  这样我们就能分辨“设 $x=\str{a}$”
  与
  “设 $x=a$”。
  符号不是变量\Dash 它们不持有某个值，
  它们本身就是一个值。}
%</def:Alphabet>
%<*def:String>
字母表上的一个 \definend{串（string）}\index{string}
或 \definend{字（word）}\index{word over an alphabet|see {string}}
是该字母表元素的一个有限序列。
没有元素的串是
\citetext{\definend{空串（empty string）}，记作~$\emptystring$}{或许
  $\emptystring$ 是 ‘empty’ 的缩写。
  有些作者使用 $\lambda$，可能源自德语的“空”字 \textit{leer}。
  也可能只是某人正好需要个符号就用了它；传闻当被问及为何在他的 $\lambda$ 演算中使用 $\lambda$ 符号时，
  Church 回答说，“eenie, meenie, meinie, mo”
  \autocite{csse:Siren}；亦可参阅
  \url{https://www.youtube.com/watch?v=juXwu0Nqc3I}}。\index{string!empty}\index{empty string@empty string, $\emptystring$}
%</def:String>

%<*discussion:SymbolsMayContainMoreThanOneChar>
关于符号，有一点可能出乎意料：
它可以包含不止一个字符。
例如，一门编程语言
可能将 \str{if} 作为一个不可分解为单独字母的符号。
%</discussion:SymbolsMayContainMoreThanOneChar>
另一个例子是，Racket 的字母表
包含符号 \str{or} 和 \str{car}，同时
也允许像 \str{lastname} 这样的符号作为变量名。
Racket 中一个串的例子是 $\str{(or a ready)}$，它是
五个字母表元素的序列，
$\sequence{\str{(}, \str{or}, \str{a}, \str{ready},\str{)}}$。

%<*discussion:AlphSymbol>
传统上用大写希腊字母~$\Sigma$ 来表示字母表。\footnote{%
  封面内页列出了希腊字母及其读音。}
我们用小写希腊字母来命名串
（不过我们把~$\TMfcn$ 留作它用），
并用相应的小写罗马字母来表示串中的符号，如
$\sigma=\sequence{s_0,\ldots\, s_{n-1}}$
和 $\tau=\sequence{t_0,\ldots\, t_{m-1}}$。
%</discussion:AlphSymbol>
\label{def:StringLength}
%<*def:StringLength>
串 $\sigma$ 的 \definend{长度（length）}，
记作 \definend{$\lh{\sigma}$}，\index{string!length}\index{length!of a string}
是它所包含的符号个数，即~$n$。
特别地，
空串的长度是 $\lh{\emptystring}=0$。
%</def:StringLength>

有时我们用方括号 $\sigma[i]$ 来代替下标记法 $s_i$。
这样做的好处是我们可以用 $\sigma[-1]$ 表示最后一个符号，
用 $\sigma[-2]$ 表示倒数第二个符号，等等。
我们还用 $\sigma[i\mathord{:} j]$ 表示
位于~$s_i$ 和~$s_j$ 之间的符号子序列，
包含~$s_i$ 但不包含~$s_j$；
用 $\sigma[i\mathord{:}]$ 表示从~$s_i$ 开始的尾部；
用 $\sigma[\mathord{:} j]$ 表示 $\sigma[0\mathord{:} j]$。

%<*dis:OmitCommas>
为方便起见，我们通常使用
符号为单字符的字母表，这样就能
省略括号和逗号来书写串。
也就是说，比起 $\sequence{\str{a},\str{b},\str{c}}$，
我们更喜欢 $\sigma=\str{abc}$。\footnote{%
  要理解为什么当我们去掉逗号时，希望字母表由
  单字符符号组成，可以考虑
  $\Sigma=\set{\str{a}, \str{aa}}$
  和串 $\str{aaa}$。
  没有逗号时，这个串是有歧义的，它可能
  表示 ~$\sequence{\str{a},\str{a}, \str{a}}$、$\sequence{\str{a},\str{aa}}$、
  $\sequence{\str{aa},\str{a}}$，
  或~$\sequence{\str{aaa}}$。}
%</dis:OmitCommas>
这样做的
缺点在于，没有了尖括号，
空串就什么都不是了，因此
我们专门保留了一个单独的符号~$\emptystring$。\footnote{%
  省略
  括号和逗号也会模糊符号
  和单符号串之间的区别，即模糊了 \str{a} 和 $\sequence{\str{a}}$ 的区别。
  然而，这样做实在太方便了，我们愿意接受这个
  缺点。}

%<*def:Bitstrings>
由比特字符组成的字母表是
\definend{$\B=\set{\str{0},\str{1}}$}
（有时我们也用 $\B$ 表示比特本身的集合 $\set{0,1}$）。
$\B$ 上的串称为 \definend{比特串（bitstrings）}，\index{bitstring}\index{bitstring!set of@set of, $\B$}
或 \definend{比特串（bit strings）}。\index{bit string|see{bitstring}}
%</def:Bitstrings>

%<*def:KleeneStar>
设~$\Sigma$ 是一个字母表，
对于 $k\in\N$，
该字母表上长度为~$k$ 的串的集合是
\definend{$\Sigma^k\!$}。
$\Sigma$ 上所有有限长度串的集合是
\definend{$\kleenestar{\Sigma}=\cup_{k\in\N}\Sigma^k\!$}。
星号符号称为 \definend{Kleene 星（Kleene star）}，\index{Kleene star}
读作“star”。
%</def:KleeneStar>

%<*def:StringOps>
串的操作只有少数几种。
设 $\sigma=\sequence{s_0,\ldots\, s_{n-1}}$ 和
$\tau=\sequence{t_0,\ldots\, t_{m-1}}$ 是同一字母表上的两个串。
\definend{拼接 $\sigma\concat\tau$} 或 \definend{$\sigma\tau$}\index{concatenation!of strings}\index{string!concatenation}
是指将第二个串附加到第一个串之后，$\sigma\concat\tau=\sequence{s_0,\ldots\, s_{n-1},t_0,\ldots\, t_{m-1}}$。
\label{def:substring}
若
$\sigma=\tau_0\concat \cdots \concat\tau_{k-1}$，我们就称 $\sigma$ \definend{分解}\index{string!decomposition}
为这些 $\tau$，并且每个 $\tau_i$ 都是 $\sigma$ 的一个
\definend{子串}\index{substring}\index{string!substring}。
第一个子串 $\tau_0$ 是 $\sigma$ 的一个
\definend{前缀}\index{prefix of a string}\index{string!prefix}。
最后一个子串 $\tau_{k-1}$ 是一个
\definend{后缀}。\index{suffix of a string}\index{string!suffix}

一个串的
\definend{幂}\index{power of a string}\index{string!power}
或
\definend{重复}\index{replication of a string}\index{string!replication}
是指它自身迭代拼接的结果，因此 $\sigma^2=\sigma\concat\sigma$，
$\sigma^3=\sigma\concat\sigma\concat\sigma$，等等。
我们规定 $\sigma^1=\sigma$，且 $\sigma^0=\emptystring$。
串的
\citetext{\definend{反转 $\reversal{\sigma}$}}{%
  当前最实用的字符串概念——Unicode 标准——并没有定义字符串反转操作。
  所有反转字符串的朴素方法在处理任意 Unicode 字符串时都会遇到问题，因为这些字符串可能包含非 ASCII
  字符、组合字符、连字、多种语言的双向文本等等。
  例如，仅仅反转字符串中的字符（Unicode 标量值）就可能导致组合附加符号附着到错误的字符上。
  另一个例子是：如何反转 \str{ab<backspace>ab}？
  Unicode 联盟并没有费力去定义字符串的反转，因为现实中并没有这个需求。
  （来自 \url{https://qntm.org/trick}。）}\index{reversal!of a string}\index{string!reversal}
是指将符号按相反顺序排列：$\reversal{\sigma}=\sequence{s_{n-1}, \ldots\, s_0}$。
空串的反转是
$\reversal{\emptystring}=\emptystring$。
%</def:StringOps>

例如，
设 $\Sigma=\set{\str{a},\str{b},\str{c}}$，并令
$\sigma=\str{abc}$，$\tau=\str{bbaac}$。
那么拼接 $\sigma\tau$ 就是 \str{abcbbaac}。
三次幂 $\sigma^3$ 是 \str{abcabcabc}，
而反转 $\reversal{\tau}$ 是 \str{caabb}。
%<*def:Palindrome>
\definend{回文}\index{palindrome}
是等于其自身反转的串。
例如 $\alpha=\str{abba}$、$\beta=\str{cdc}$ 和 $\emptystring$。
%</def:Palindrome>

\begin{exercises}


\begin{exercise}
设 $\sigma=\str{10110}$ 和 $\tau=\str{110111}$ 为比特串。
求下列各项。
\begin{exparts*}
\item $\sigma\concat \tau$
\item $\sigma\concat\tau\concat\sigma$
\item $\reversal{\sigma}$
\item $\sigma^3$
\item $\str{0}^3\concat\sigma$
\end{exparts*}
\begin{answer}
\begin{exparts}
\item $\sigma\concat \tau=\str{10110}\concat\str{110111}=\str{10110110111}$。
\item $\sigma\concat\tau\concat\sigma=
   \str{10110}\concat\str{110111}\concat\str{10110}=\str{1011011011110110}$
\item $\reversal{\sigma}=\reversal{\str{10110}}=\str{01101}$
\item $\sigma^3=\str{10110}\concat\str{10110}\concat\str{10110}
               =\str{101101011010110}$
\item $\str{0}^3\concat\sigma=\str{000}\concat\str{10110}=\str{00010110}$
\end{exparts}
\end{answer}
\end{exercise}

\begin{exercise}
设字母表为 $\Sigma=\set{\str{a},\,\str{b},\,\str{c}}$。
假设 $\sigma=\str{ab}$ 且 $\tau=\str{bca}$。
求下列各项。
\begin{exparts*}
\item $\sigma\concat\tau$
\item $\sigma^2\concat\tau^2$
\item $\reversal{\sigma}\concat\reversal{\tau}$
\item $\sigma^3$
\end{exparts*}
\begin{answer}
\begin{exparts}
\item \str{abbca}
\item \str{ababbcabca}
\item \str{baacb}
\item \str{ababab}
\end{exparts}
\end{answer}
\end{exercise}

\begin{exercise}
设
$\lang=\set{\sigma\in\kleenestar{\B}\suchthat
   \text{$\lh{\sigma}=4$ 且 $\sigma$ 以 \str{0} 开头}}$。
该语言中有多少个元素？
\begin{answer}
长度为 $4$ 的比特串共有 $2^4=16$ 个，其中一半以 \str{0} 开头，因此该语言中有 $8$ 个串。
\end{answer}
\end{exercise}

\begin{exercise}
假设 $\Sigma=\set{\str{a},\,\str{b},\str{c}}$
且 $\sigma=\str{abcbccbba}$。
\begin{exparts*}
\item \str{abcb} 是 $\sigma$ 的前缀吗？
\item \str{ba} 是后缀吗？
\item \str{bab} 是子串吗？
\item $\emptystring$ 是后缀吗？
\end{exparts*}
\begin{answer}
\begin{exparts}
\item 是。
\item 是。
\item 否。
\item 是。
\end{exparts}
\end{answer}
\end{exercise}

\begin{exercise}
$\lh{\sigma}$、$\lh{\tau}$ 和 $\lh{\sigma\concat\tau}$ 之间有什么关系？
你必须说明理由。
\begin{answer}
拼接的长度等于两个长度之和：
$\lh{\sigma\concat\tau}=
\lh{\sequence{s_0,\ldots\, s_{i-1},t_0,\ldots\, t_{j-1}}}=i+j=\lh{\sigma}+\lh{\tau}$。
\end{answer}
\end{exercise}

\begin{exercise}
串的拼接操作遵循简单的代数规律。
对于下列每一条，判断其真假。
若为真，则予以证明；若为假，则举出反例。
\begin{exparts*}
  \item $\alpha\concat\emptystring=\alpha$
    且 $\emptystring\concat\alpha=\alpha$
  \item $\alpha\concat\beta=\beta\concat\alpha$
  \item $\reversal{\alpha\concat\beta}=\reversal{\beta}\concat\reversal{\alpha}$
  \item $\reversal{\reversal{\alpha}}=\alpha$
  \item $\reversal{\alpha^i}=\alpha^i$
\end{exparts*}
\begin{answer}
\begin{exparts}
\item 真。
  第一式的拼接结果为
  $\alpha\concat\emptystring=\sequence{a_0,\ldots\, a_{i-1}}\concat\sequence{}=\sequence{a_0,\ldots\, a_{i-1}}=\alpha$。
  第二式的证明方法类似。
\item 假。
  反例为比特串 $\alpha=\str{0}$ 和 $\beta=\str{1}$。
\item 假。
  与前一项类似，
  反例为比特串 $\alpha=\str{0}$ 和 $\beta=\str{1}$。
\item 真。
  当 $\alpha=\sequence{a_0,\ldots\, a_{i-1}}$
  时，我们有 $\reversal{\alpha}=\sequence{a_{i-1},\ldots\, a_{0}}$，
  因此 $\reversal{\reversal{\alpha}}=\sequence{a_{0},\ldots\, a_{i-1}}$。
\item 假。
  反例为取幂次 $i=1$ 的比特串 $\alpha=\str{01}$。
\end{exparts}
\end{answer}
\end{exercise}

% \begin{exercise}
% Show that string concatenation is not commutative, that there are
% strings $\sigma$ and $\tau$ so that
% $\sigma\concat\tau\neq\tau\concat\sigma$.
% \begin{answer}
% One example using $\B$ is $\sigma=\sequence{0}$ and $\tau=\sequence{1}$.
% Then $\sigma\concat\tau=\sequence{0,1}$ while
% $\tau\concat\sigma=\sequence{1,0}$.
% \end{answer}
% \end{exercise}

% \begin{exercise}
% In defining decomposition above we have
% `$\sigma=\tau_0\concat \cdots \concat\tau_{n-1}$',
% without parentheses on the right side.
% This takes for granted that
% the concatenation operation is associative,
% that no matter how we parenthesize it we get the same string.
% Prove this.
% \hint{} use induction on the number of substrings, $n$.
% \end{exercise}

% \begin{exercise}
% Prove that this constructive definition of string power is equivalent
% to the one above.
% \begin{equation*}
%   \sigma^n =
%     \begin{cases}
%       \emptystring  &\case{if $n=0$}  \\
%        \sigma^{n-1}\concat\sigma &\case{if $n>0$}
%     \end{cases}
% \end{equation*}
% \end{exercise}

% \begin{exercise} \index{Levi's lemma}
% (Levi's lemma)
% Prove that for all strings $\alpha,\beta,\gamma,\delta\in\kleenestar{\Sigma}$,
% if $\alpha\concat\beta=\gamma\concat\delta$
% then there exists a string $\sigma$ such that either
% (1)~$\alpha\concat\sigma=\gamma$
% or (2)~$\alpha=\gamma\concat\sigma$.
% That is, there is a string in the middle
% that can be grouped to one side or the other.
% \hint{} case~(1) applies when $\lh{\alpha}\leq\lh{\gamma}$.
% \end{exercise}
\end{exercises}
\index{string|)}

\endinput